\section{Conclusion}\label{Conclusion}
\pagestyle{mypagestyle}
\begin{spacing}{1.75}
	
	\MUFGfont{
		\subsection{Validation conclusion}
		This validation assessed the FX simulation model used within the PFE framework, focusing on:
		(i) the FX simulation specification (Model document 1.2, Section 2.4) and
		(ii) the FX calibration approach (Model document 1.1, Section 3.3).
		
		The model simulates FX rates (all currencies versus USD) using a GBM construction:
		a driftless Brownian driver with a lognormal transform, where the deterministic mean function is derived from discount factors
		(foreign vs USD), implying a forward-consistent drift. Volatility is calibrated from historical 5-day log returns using a three-year history.
	}
	
	\MUFGfont{
		\subsection{Assessment by validation dimension}
		\subsubsection{Conceptual soundness}
		The model is conceptually consistent with a standard PFE exposure simulation approach:
		\begin{itemize}
			\item GBM dynamics are mathematically well-defined and computationally robust;
			\item forward-implied drift via discount factors is consistent with covered interest parity and arbitrage-free relationships at $t=0$;
			\item the USD reporting convention and derivation of cross rates reduce dimensionality and enforce internal consistency.
		\end{itemize}
		However, the model does not capture higher-order FX dynamics (smile/skew, stochastic volatility, jumps), which may impact tail exposures
		in stress regimes and requires explicit limitation statements and compensating controls.
	}
	
	\MUFGfont{
		\subsubsection{Calibration methodology}
		The calibration approach is straightforward and operationally implementable. The key challenge is that volatility is estimated from a fixed historical
		window using 5-day returns, which may be sensitive to:
		(i) regime changes,
		(ii) the chosen sampling frequency,
		(iii) missing data and outlier handling.
		Validation therefore requires: (a) stability evidence across windows and frequencies and (b) clear override governance.
	}
	
	\MUFGfont{
		\subsubsection{Data integrity and implementation}
		Data quality is critical: FX spot history must be consistent in sourcing and time-of-day, and curve identifiers used in the mean function must align
		with the valuation curve build as-of date.
		Implementation should include unit tests verifying: moment consistency ($\mathbb{E}[Y(t)]\approx g(t)$), variance scaling,
		and cross-rate identities derived from simulated USD legs.
	}
	
	\MUFGfont{
		\subsubsection{Performance and monitoring}
		Given the model's intended use for PFE, performance validation should emphasize:
		\begin{itemize}
			\item distributional tests on standardized residuals / PIT;
			\item exception-based conservativeness tests on key FX quantiles;
			\item calibration stability and break detection;
			\item PFE sensitivity to volatility stresses and configuration choices.
		\end{itemize}
	}
	
	\MUFGfont{
		\subsection{Overall model risk rating and decision}
		\textbf{Provisional rating: \underline{Amber}}. The model is suitable for PFE exposure simulation provided that the following conditions are met:
		\begin{enumerate}
			\item completion of the performance and stability tests defined in Section \ref{Performance}, including sensitivity analyses on calibration
			window length and return frequency;
			\item implementation and evidence of unit tests and cross-rate consistency checks (Section \ref{Implementation});
			\item formalization of calibration governance, overrides, and monitoring KPIs with escalation thresholds (Section \ref{Governance});
			\item explicit documentation of model limitations and permitted-use statements (Section \ref{AssumptionsLimitations}).
		\end{enumerate}
	}
	
	\MUFGfont{
		\subsection{Summary conclusion table}
		\begin{table}[!h]
			\centering
			\arrayrulecolor{mufggrey}
			\caption{\MUFGfont{Conclusion Summary for \ModelID}}
			\vspace{5pt}
			\begin{tabular}{|p{5.0cm}|p{3.0cm}|p{8.0cm}|}
				\hline
				\cellcolor{mufglightgrey}\MUFGfont{\textbf{Dimension}} &
				\cellcolor{mufglightgrey}\MUFGfont{\textbf{Assessment}} &
				\cellcolor{mufglightgrey}\MUFGfont{\textbf{Key Comments / Conditions}} \\ \hline
				
				\MUFGfont{Conceptual Soundness} &
				\MUFGfont{Satisfactory} &
				\MUFGfont{Forward-implied drift and GBM structure are appropriate for PFE; limitations on tails/smile require explicit disclosures and controls.} \\ \hline
				
				\MUFGfont{Calibration Methodology} &
				\MUFGfont{Satisfactory (conditional)} &
				\MUFGfont{Historical vol from 5-day returns requires stability evidence and sensitivity testing; define override governance.} \\ \hline
				
				\MUFGfont{Data Integrity} &
				\MUFGfont{Satisfactory (conditional)} &
				\MUFGfont{Require completeness/outlier checks on FX history and curve alignment controls; auditability of calibration snapshots.} \\ \hline
				
				\MUFGfont{Implementation} &
				\MUFGfont{Satisfactory (conditional)} &
				\MUFGfont{Must demonstrate unit tests: $\mathbb{E}[Y(t)]\!=\!g(t)$, variance scaling, and cross-rate identity checks.} \\ \hline
				
				\MUFGfont{Performance \& Monitoring} &
				\MUFGfont{Requires enhancement} &
				\MUFGfont{Implement PIT/exception tests, calibration break detection, and stress overlays; define escalation thresholds.} \\ \hline
				
			\end{tabular}
		\end{table}
	}
	
	\MUFGfont{
		\subsection{Action plan and next review}
		An action plan should be agreed between Model Owner and Model Risk to address conditional items (tests, monitoring, governance documentation),
		with target completion dates. Following completion, the model risk rating should be reassessed and the model moved to \textbf{Green} if evidence is satisfactory.
		Periodic review requirements are set out in Section \ref{Governance}.
	}
	
\end{spacing}
\clearpage
